% =====================================================================
%  Data Science in Finance and Economics (DSFE) -- AIMS Press
%  Special Issue: Machine Learning in Economics and Finance
%
%  Manuscript draft. Formatting follows the DSFE / AIMS Press author
%  instructions: Times New Roman 12pt, 15pt leading, author-year
%  citations, three-line (booktabs) tables with captions ABOVE, figure
%  captions BELOW, sentence-case title and headings.
%
%  Self-contained: inline thebibliography, compiles in one pdflatex pass.
%  Build:  pdflatex main  ->  pdflatex main   (twice for cross-refs)
% =====================================================================
\documentclass[12pt,a4paper]{article}

% Fonts: Times New Roman equivalent
\usepackage[T1]{fontenc}
\usepackage[utf8]{inputenc}
\usepackage{newtxtext,newtxmath}

% Language: British English
\usepackage[british]{babel}

% Page geometry and spacing
\usepackage[a4paper,margin=2.5cm]{geometry}
\usepackage{setspace}
\setstretch{1.25}

% Maths, tables, graphics
\usepackage{amsmath}
\usepackage{booktabs}
\usepackage{graphicx}
\graphicspath{{figures/}}
\usepackage{array}
\usepackage{caption}
\captionsetup{font=small,labelfont=bf,labelsep=period}
\captionsetup[table]{skip=4pt}

% Authors / affiliations
\usepackage{authblk}
\renewcommand\Authfont{\normalsize}
\renewcommand\Affilfont{\small\itshape}
\setlength{\affilsep}{0.4em}

% Citations: author-year (DSFE / AIMS style)
\usepackage[round,authoryear,sort&compress]{natbib}
\setcitestyle{aysep={,}}
\bibpunct{(}{)}{;}{a}{,}{,}

% Section heading style: 12pt, sentence case
\usepackage{titlesec}
\titleformat{\section}{\normalsize\bfseries}{\thesection.}{0.5em}{}
\titleformat{\subsection}{\normalsize\bfseries}{\thesubsection.}{0.5em}{}
\titleformat{\subsubsection}{\normalsize\itshape}{\thesubsubsection.}{0.5em}{}
\titlespacing*{\section}{0pt}{1.4ex plus .2ex}{0.8ex}
\titlespacing*{\subsection}{0pt}{1.2ex plus .2ex}{0.6ex}

% Hyperlinks (load last)
\usepackage[colorlinks=true,linkcolor=black,citecolor=black,urlcolor=blue!60!black]{hyperref}
\urlstyle{same}  % render DOIs in body font, not typewriter

\begin{document}

% =====================================================================
%  TITLE BLOCK
% =====================================================================
\title{\large\bfseries Diagnosing and mitigating the neutral-class deficit in financial
sentiment analysis: An explainability-guided approach on multi-source investor text}

\author[1,*]{Dominic Hubble}
\affil[1]{Department of Computer Science, Loughborough University, Loughborough, LE11 3TU, UK}

\date{}
\maketitle

\renewcommand{\thefootnote}{\fnsymbol{footnote}}
\footnotetext[1]{\textbf{*Correspondence:} Email: dominichubble@gmail.com;
Tel: +44 (0)1509 222222; Postal address: Department of Computer Science,
Loughborough University, Loughborough, LE11 3TU, United Kingdom.}
\renewcommand{\thefootnote}{\arabic{footnote}}

% =====================================================================
%  ABSTRACT
% =====================================================================
\begin{abstract}
\noindent
Domain-adaptive transformer models such as FinBERT are the de facto standard for
financial sentiment classification, yet their well-documented weakness on the
neutral class is typically reported as an aggregate error rate rather than
explained or repaired. This study treats that weakness as a diagnosable,
addressable phenomenon. We assemble a 200-sample evaluation benchmark that
deliberately reflects the noisy, multi-source distribution of deployed financial
text, spanning retail discussion (Reddit), microblog reactions (X/Twitter), and
professional news, rather than the clean, professionally annotated sentences of
existing corpora. Benchmark labels are validated by a large-language-model
cross-annotation pass (Cohen's $\kappa = 0.99$ against the author) and an
independent human annotator on a stratified subset (Cohen's $\kappa = 0.90$). On
this benchmark, FinBERT attains 82.0\% accuracy and 80.3\% macro F1, outperforming
a finance-lexicon keyword baseline by 13.5 and 12.4 percentage points; the gap is
significant under a paired McNemar test ($p = 0.002$) and confirmed by
non-overlapping stratified-bootstrap confidence intervals. We then use a local
interpretability method (LIME) not as a presentation layer but as a diagnostic
instrument, showing that FinBERT's 53.3\% neutral recall is driven by a single
systematic failure mode: contextually neutral risk vocabulary (e.g.\ ``cautious'',
``inflation'') is consistently attributed to the negative class. Motivated by this
diagnosis, a targeted fine-tuning intervention on a small (300-sample), disjoint
corpus of neutral-but-lexically-loaded text recovers neutral recall to 71.7\%
($+18.4$ points) and lifts overall accuracy to 86.5\%, with negligible degradation
on the polarised classes. The contribution is methodological: an
explainability-guided error-analysis-and-mitigation loop for financial sentiment
models, validated on realistic multi-source text with appropriate paired
significance testing. We report supporting analyses (a finance-emotion enrichment
layer and a sentiment-price correlation engine) transparently, including their
negative and null results.
\end{abstract}

\noindent\textbf{Keywords:} financial sentiment analysis; FinBERT; explainable AI;
LIME; neutral-class recall; domain adaptation; targeted fine-tuning; investor
sentiment; multi-source text; model error analysis\\[2pt]
\noindent\textbf{JEL Codes:} C45; C52; C53; G14; G17; G41

\vspace{0.5em}
\noindent\textbf{Abbreviations:} XAI, explainable artificial intelligence; LIME,
local interpretable model-agnostic explanations; SHAP, Shapley additive
explanations; LLM, large language model; API, application programming interface.

% =====================================================================
\section{Introduction}
% =====================================================================
Investor sentiment is a recognised driver of asset prices, trading volume and
volatility, capable of producing temporary mispricings that persist until
information is absorbed \citep{tetlock2007,baker2006}. The migration of investor
discourse onto social media and online communities has made it possible to
quantify this sentiment at scale, and domain-adaptive transformer models, chiefly
FinBERT \citep{araci2019}, have become the standard tool for doing so, capturing
implicit financial phrasing (``guidance raised'', ``missed expectations'', ``strong
headwinds'') that earlier lexicon and bag-of-words methods cannot.

A recurring and well-documented limitation of FinBERT is its weak performance on
the \emph{neutral} class: the model is biased toward the polarised classes on
which its pre-training corpus is richest, and routinely mislabels factual,
sentiment-free statements as positive or negative. This limitation is almost
always reported as a number (a low neutral recall or a confusion-matrix row) and
left there. What is missing is, first, a \emph{mechanistic} account of \emph{why}
the model fails on neutral text, and second, evidence that the failure is
\emph{addressable} with a proportionate intervention. Existing surveys of
explainable AI (XAI) in finance note a parallel gap: interpretability methods are
overwhelmingly used either for post-hoc presentation or for offline auditing,
rather than as instruments that feed back into model improvement
\citep{yeo2025}.

This paper closes that loop. Our central claim is that FinBERT's neutral-class
deficit is not diffuse noise but a \emph{specific, lexically-localised, and
repairable} failure mode, and that a local interpretability method can be used to
diagnose it precisely enough to motivate a targeted, measurable fix. Concretely,
we make the following contributions:
\begin{enumerate}
\item \textbf{A realistic multi-source benchmark and a rigorous baseline
comparison.} We construct a 200-sample financial-sentiment benchmark reflecting
the true mixture of retail, microblog and news text that a deployed system
ingests, and we validate its labels with both LLM cross-annotation and an
independent human annotator. We compare FinBERT against a finance-lexicon keyword
baseline using a \emph{paired} significance test (McNemar) and stratified-bootstrap
confidence intervals, a level of statistical rigour that comparable applied studies
frequently omit.
\item \textbf{An explainability-guided diagnosis of the neutral-class deficit.}
Using LIME as a diagnostic instrument rather than a user-facing feature, we show
that FinBERT's neutral failures are driven by a consistent mechanism: contextually
neutral risk vocabulary is systematically attributed to the negative class.
\item \textbf{A targeted, evidence-led mitigation.} Guided by the diagnosis, we
fine-tune on a small, disjoint corpus of neutral-but-lexically-loaded financial
text and recover neutral recall by 18.4 percentage points while preserving
polarised-class performance, demonstrating that the diagnosed failure is
proportionately addressable.
\end{enumerate}

These results are obtained within, and reported alongside, a complete deployed
analytics system; the system's additional components (a finance-emotion enrichment
layer, a lightweight aspect extractor, and a sentiment-price correlation engine)
are described as engineering context and their evaluations, including negative and
null findings, are reported transparently rather than promoted to headline claims.
The remainder of the paper is organised as follows. Section~\ref{sec:related}
reviews related work; Section~\ref{sec:data} describes the data;
Section~\ref{sec:method} sets out the methodology; Section~\ref{sec:results}
presents results; Section~\ref{sec:discussion} discusses implications and threats
to validity; and Section~\ref{sec:conclusion} concludes.

% =====================================================================
\section{Related work}\label{sec:related}
% =====================================================================
\subsection{Sentiment, text and markets}
Behavioural finance establishes that sentiment and psychological factors produce
temporary, economically meaningful mispricings \citep{baker2006}. \citet{tetlock2007}
showed that negative language in financial news predicts short-term downward
pressure on returns with subsequent reversion; \citet{antweiler2004} had already
established that user-generated message-board text carries statistically
significant information about volatility, and \citet{schumaker2009} reported
comparable short-horizon predictive content in breaking financial news.
\citet{bollen2011} extended this to social
media at scale, deriving daily mood indices from Twitter that predicted Dow Jones
movements, while \citet{da2015} reached a similar conclusion from aggregate search
volumes, and \citet{garcia2013} showed news-sentiment effects are strongest during
recessions. More recent work documents the distinct behaviour of retail
communities: \citet{long2023} show that Reddit sentiment during the GameStop
episode diverged sharply from institutional news, and cryptocurrency markets
exhibit comparable retail-driven sensitivity to Twitter sentiment
\citep{kraaijeveld2020}. The collective implication is that single-source sentiment
signals are structurally incomplete, motivating multi-source measurement.

\subsection{From lexicons to domain-adaptive transformers}
Early financial sentiment systems relied on domain-specific lexicons;
\citet{loughran2011} demonstrated that general-purpose word lists misclassify
finance vocabulary and produced the finance-specific dictionary that remains the
standard lexicon baseline. Lexicon methods are transparent and fast but fail on
negation, implicit polarity and context-dependent phrasing. The transformer
architecture \citep{vaswani2017} and bidirectional pre-training \citep{devlin2019}
changed the field, and FinBERT \citep{araci2019} adapted this to finance by
pre-training on financial corpora, consistently outperforming general models and
lexicons on the Financial PhraseBank benchmark \citep{malo2014}. Parallel FinBERT
variants target financial communications and filings
\citep{yang2020,huang2023}.

Since 2023, large language models (LLMs) have been applied to financial sentiment
with mixed but informative results. ChatGPT-based scores have been found to
outperform FinBERT on some benchmarks and to carry incremental return-predictive
information \citep{fatouros2023,lopezlira2023}, while other work shows LLM sentiment
is brittle under adversarial rephrasing \citep{leippold2023} and unreliable on
informal Reddit-register text \citep{deng2023}. Domain LLMs such as BloombergGPT
\citep{wu2023} demonstrate strong zero-shot performance but are not publicly
available. The consensus is that LLMs have narrowed the per-sentence accuracy gap
with FinBERT on clean benchmarks but have not displaced it where latency, cost and
reproducibility are first-order constraints, which is the setting of this study.
Across this literature, the neutral class is repeatedly identified as the hardest
and least-studied category, but is rarely diagnosed or remediated directly.

\subsection{Explainable AI in finance: presentation versus diagnosis}
\citet{doshivelez2017} frame interpretability as explaining model behaviour in
human-understandable terms wherever objectives cannot be fully specified in
advance, which is almost always the case in finance. LIME \citep{ribeiro2016}
explains individual predictions by fitting a sparse linear surrogate over
perturbations of the input, identifying the tokens most responsible for a
prediction; SHAP \citep{lundberg2017} provides game-theoretic guarantees at higher
computational cost. Surveys find that XAI in finance is dominated by post-hoc LIME
and SHAP but that deployment within interactive tools, and use as a feedback
mechanism, remain rare \citep{yeo2025}. \citet{sravankumar2025} integrate XAI with
deep learning for sentiment-driven stock prediction. Two cautionary
results bound the claims that interpretability can support: perturbation-based
explanations can be adversarially fooled \citep{slack2020}, and intrinsically
interpretable models may be preferable in high-stakes settings \citep{rudin2019}.
We therefore use LIME strictly as a \emph{diagnostic aid}, a hypothesis-generating
instrument whose output is corroborated against the confusion matrix and an
intervention, rather than as a guarantee of correctness. The gap this paper
addresses is thus narrow and specific: the use of local interpretability to
\emph{diagnose a known model failure mode precisely enough to drive a targeted,
validated mitigation}, evaluated on realistic multi-source financial text.

% =====================================================================
\section{Data}\label{sec:data}
% =====================================================================
\subsection{Multi-source corpus}
The deployed system ingests text from three complementary channels via documented
official interfaces: Reddit (long-form discussion from finance subreddits, via
PRAW), X/Twitter (high-frequency reactions, via the official X API v2), and
financial news (professionally sourced articles, via NewsAPI). All collection is
restricted to publicly available text; author-identifying fields (handles, display
names, geolocation) are stripped at ingestion, so that each stored record contains
only sanitised text, a timestamp, a source identifier, and any resolved ticker
mentions. Duplicate and near-duplicate detection (content-hashing and TF-IDF
cosine similarity) and minimum-length thresholds are applied. At the time of
analysis the corpus comprised more than 150{,}000 labelled records. We note a
material data-layer limitation up front: tightened X API v2 quotas during
2023 to 2024 bounded the Twitter subcorpus well below the Reddit and news
subcorpora, so the Twitter signal is treated as secondary and corroborating
throughout (see Section~\ref{sec:discussion}). The project received a favourable
ethical opinion from the Loughborough University Ethics Review Sub-Committee through
the University's ethical review system; full approval and consent details are given
in the Ethics statement.

\subsection{Evaluation benchmark and label validation}\label{sec:benchmark}
Existing benchmarks such as Financial PhraseBank \citep{malo2014} consist of
professionally written news sentences and do not reflect the informal, noisy
language the deployed system processes. We therefore constructed a purpose-built
200-sample benchmark spanning the three ingestion channels, balanced across
classes: 70 positive (earnings beats, guidance raises, upgrades), 70 negative
(misses, downgrades, sell-offs), and 60 neutral (factual announcements, regulatory
disclosures, on-par results). Primary ground-truth labels were assigned by the
author, cross-referencing each example against the Financial PhraseBank taxonomy
for consistency.

Because single-annotator labelling is a threat to validity, we validated the labels
in two independent ways. First, a large language model, given only the 200 raw
texts with no access to the author's labels and instructed to apply the Financial
PhraseBank taxonomy, produced a second label set; agreement was 199/200 (Cohen's
$\kappa = 0.9925$, SE $= 0.0075$) \citep{cohen1960}. We interpret this conservatively
as an \emph{internal-consistency} check rather than a human inter-rater coefficient,
because the model and the author share overlapping training distributions and the
figure will overstate a human study. Second, an independent human annotator, a
finance undergraduate with equity-research internship experience and no involvement
in the project, labelled a 100-sample stratified subset (34/33/33) under the same
instructions. Author-to-human agreement was 93.0\% (Cohen's $\kappa = 0.895$);
three-way Fleiss' $\kappa$ \citep{fleiss1971} across the author, the human and the
second label set on the subset was 0.893. Both coefficients fall in the ``almost
perfect'' band \citep{landis1977}, and the seven disagreements between the human and
the author clustered on the neutral/positive boundary, consistent with the failure
mode analysed in Section~\ref{sec:results}. These checks support, but do not fully
substitute for, a multi-annotator panel (Section~\ref{sec:discussion}).

\subsection{Fine-tuning corpus}\label{sec:ftcorpus}
For the mitigation experiment (Section~\ref{sec:ftmethod}), we assembled a
300-sample corpus of financial text that is sentiment-neutral but lexically
sentiment-loaded, drawn from regulatory filings (8-K, 10-Q excerpts), central-bank
statements, and earnings-release language of the kind that produces FinBERT's
residual neutral confusions. This corpus was held strictly disjoint from the
evaluation benchmark (no text overlap; a source-date partition was enforced).

% =====================================================================
\section{Methodology}\label{sec:method}
% =====================================================================
\subsection{Classifier and baseline}
The primary classifier is FinBERT, loaded from the public \texttt{ProsusAI/finbert}
checkpoint and used in zero-shot inference (no task-specific fine-tuning in the main
comparison). It outputs a three-class label (positive/neutral/negative) with a
softmax probability distribution, which is retained in full for the downstream
enrichment described in Section~\ref{sec:supporting}. The baseline is a
finance-lexicon keyword classifier built on domain-specific vocabulary,
representing the best-practice lexicon approach \citep[cf.][]{loughran2011} and
providing both a meaningful lower anchor and a fast fallback. Text from all sources
passes through a shared preprocessing pipeline tuned for transformer input: case
preservation, URL/markup removal, stopword retention, no lemmatisation, preservation
of financial punctuation, and negation rewriting.

\subsection{Evaluation metrics}
We report accuracy; macro precision, recall and F1 (equal weight per class,
appropriate to the near-balanced distribution); per-class precision, recall and F1;
and full confusion matrices. For a class $c$ with precision $P_c$ and recall $R_c$,
the macro F1 is
\begin{equation}
\text{macro-F1} = \frac{1}{|\mathcal{C}|}\sum_{c\in\mathcal{C}}
\frac{2 P_c R_c}{P_c + R_c}, \qquad \mathcal{C}=\{\text{pos},\text{neg},\text{neu}\}.
\end{equation}
Inference time is recorded to contextualise deployment trade-offs.

\subsection{Paired significance testing}\label{sec:sigtest}
Comparing two classifiers on the same test set by accuracy alone cannot separate a
real performance gap from sampling variation. We therefore apply McNemar's test on
discordant pairs, the samples where exactly one classifier is correct
\citep{dietterich1998}. With $b$ samples correct only for the baseline and $c$
correct only for FinBERT, the continuity-corrected statistic is
\begin{equation}
\chi^2_{cc} = \frac{(|b-c|-1)^2}{b+c},
\end{equation}
and we additionally report the exact two-sided binomial $p$-value. To quantify
estimation precision at this sample size, we compute non-parametric 95\% confidence
intervals by stratified bootstrap resampling \citep{efron1993}: 1{,}000 replicates,
sampled with replacement at the per-class level to preserve class prevalence, with a
fixed seed.

\subsection{LIME as a diagnostic instrument}\label{sec:limemethod}
To explain individual predictions we apply LIME \citep{ribeiro2016}, which perturbs
the input by masking tokens, queries FinBERT on the perturbed variants, and fits a
sparse local linear surrogate whose weights identify the tokens most responsible for
the prediction. Rather than treating these explanations as a presentation layer, we
use them to test a hypothesis about \emph{why} neutral classification fails: we run
LIME on representative neutral texts that FinBERT misclassifies and inspect which
tokens drive the erroneous polarity. The random state is fixed for reproducibility
of the reported explanations. Consistent with \citet{slack2020} and \citet{rudin2019},
LIME output is treated as corroborating evidence, not proof, and is cross-checked
against the confusion matrix and the subsequent intervention.

\subsection{Targeted fine-tuning protocol}\label{sec:ftmethod}
Guided by the LIME diagnosis, we fine-tuned the base \texttt{ProsusAI/finbert} model
on the disjoint 300-sample neutral-loaded corpus (Section~\ref{sec:ftcorpus}) for
three epochs at a learning rate of $2\times10^{-5}$, batch size 16, with a 90/10
train/validation split and a fixed seed; training completed in approximately six
minutes on a single consumer GPU. The fine-tuned model was then re-evaluated against
the \emph{unchanged} 200-sample benchmark. This is deliberately framed as a small,
falsifiable test of whether the diagnosed failure is addressable, not as an attempt
to ship a new state-of-the-art classifier. A multi-seed, cross-validated extension
of this protocol is provided as reproducible code (Section~\ref{sec:reprod}).

\subsection{Supporting system components}\label{sec:supporting}
For completeness we summarise three components that constitute the deployed system
but are not the paper's primary contribution. (i)~A \emph{finance-emotion enrichment
layer} maps FinBERT's softmax vector, a normalised Shannon-entropy term, domain
lexicon matches and aspect priors onto six behavioural emotion classes (fear,
optimism, uncertainty, confidence, scepticism, mixed) in a single inference pass,
without labelled emotion data. The entropy term, which lifts the
\emph{uncertainty} score when the classifier is indecisive, is the normalised
Shannon entropy of the softmax vector $\mathbf{p}$,
\begin{equation}
H(\mathbf{p}) = -\frac{1}{\log K}\sum_{i=1}^{K} p_i \log p_i, \qquad K=3.
\end{equation}
(ii)~A \emph{lightweight aspect extractor}, a keyword-driven approximation of
aspect-based sentiment analysis \citep{pontiki2014}, matches sentences against six
financial aspect categories (revenue/earnings, growth/demand, risk/liquidity,
valuation, product/AI, macro/policy) to attach evidence snippets to each record. (iii)~A
\emph{sentiment-price correlation engine} computes Pearson and Spearman
correlation, configurable lag analysis, Granger causality \citep{granger1969} with
an augmented Dickey-Fuller stationarity pre-check \citep{dickey1979}, rolling
correlation, and market adjustment, against same-day returns. We report evaluations
of (i) and (iii) in Section~\ref{sec:results} transparently, including their null
and negative results.

% =====================================================================
\section{Results}\label{sec:results}
% =====================================================================
\subsection{Benchmark comparison and significance}
Table~\ref{tab:summary} reports the headline comparison. FinBERT outperforms the
keyword baseline by 13.5 percentage points in accuracy and 12.4 points in macro F1,
consistently across all metric dimensions. The 82.0\% accuracy, obtained on noisy
multi-source text, is broadly consistent with the $\sim$86\% reported for FinBERT on
the cleaner, high-agreement Financial PhraseBank subset \citep{araci2019,malo2014}
and arguably a more realistic measure of deployed performance; the baseline's 68.5\%
falls within the 65 to 72\% range typical of finance-lexicon methods, confirming it
is a fair anchor rather than a straw man.

\begin{table}[ht]
\centering
\caption{Classification performance of FinBERT versus the keyword baseline on the
200-sample benchmark.}
\label{tab:summary}
\begin{tabular}{lcccc}
\toprule
\textbf{Model} & \textbf{Accuracy} & \textbf{Macro precision} & \textbf{Macro recall} & \textbf{Macro F1} \\
\midrule
Keyword baseline & 68.5\% & 70.5\% & 69.1\% & 67.9\% \\
FinBERT & \textbf{82.0\%} & \textbf{85.1\%} & \textbf{80.6\%} & \textbf{80.3\%} \\
\midrule
\textit{Delta} & \textit{+13.5 pp} & \textit{+14.6 pp} & \textit{+11.5 pp} & \textit{+12.4 pp} \\
\bottomrule
\end{tabular}
\end{table}

The paired McNemar test (Table~\ref{tab:mcnemar}) rejects the null of equal error
rates: of 71 discordant pairs, FinBERT is correct in 49 versus 22 for the baseline
(exact two-sided binomial $p = 0.0018$; continuity-corrected $\chi^2 = 9.52$,
$p = 0.0020$). Stratified-bootstrap 95\% confidence intervals corroborate this:
FinBERT accuracy 82.0\% $[76.5\%, 87.0\%]$ and macro F1 80.3\% $[74.6\%, 85.4\%]$ do
not overlap the corresponding baseline intervals ($[62.5\%, 74.5\%]$ and
$[61.4\%, 73.9\%]$). The performance gap is therefore not attributable to sampling
variation on this benchmark.

\begin{table}[ht]
\centering
\caption{Paired McNemar test for the accuracy difference between FinBERT and the
keyword baseline.}
\label{tab:mcnemar}
\begin{tabular}{lc}
\toprule
\textbf{Statistic} & \textbf{Value} \\
\midrule
Discordant pairs ($n$) & 71 \\
$b$ (keyword-only correct) & 22 \\
$c$ (FinBERT-only correct) & 49 \\
Exact two-sided binomial $p$ & \textbf{0.0018} \\
$\chi^2$ (continuity-corrected, df $=1$) & 9.52 \\
Approximate $p$ from $\chi^2$ & 0.0020 \\
\bottomrule
\end{tabular}
\end{table}

\subsection{Per-class structure and the neutral deficit}
Table~\ref{tab:perclass} and the confusion matrix (Table~\ref{tab:cm}) expose where
the aggregate gap originates. FinBERT's largest gain is on the negative class, where
recall rises from 47.1\% to 94.3\%: the lexicon baseline depends on explicit negative
vocabulary and fails on implicit phrasing (``margins compressed'', ``missed by a
whisker''), whereas FinBERT's contextual encoding handles these. The critical
pattern, however, is the \emph{neutral asymmetry}: FinBERT achieves very high neutral
precision (97.0\%) but only 53.3\% recall. It is conservative, labelling text
neutral only when highly confident no sentiment signal is present, and errs toward a
polarised class whenever sentiment-bearing vocabulary appears. Of 60 neutral samples,
28 are misclassified (17 as positive, 11 as negative), almost all factual statements
containing sentiment-loaded words (e.g.\ ``the company \emph{raised} guidance to
reflect \emph{stronger} demand''). Figure~\ref{fig:bars} visualises the per-class
contrast.

\begin{table}[ht]
\centering
\caption{Per-class precision, recall and F1 on the 200-sample benchmark (support:
pos $=70$, neg $=70$, neu $=60$).}
\label{tab:perclass}
\begin{tabular}{llccc}
\toprule
\textbf{Model} & \textbf{Class} & \textbf{Precision} & \textbf{Recall} & \textbf{F1} \\
\midrule
Keyword & Positive & 80.9\% & 78.6\% & 79.7\% \\
        & Negative & 75.0\% & 47.1\% & 57.9\% \\
        & Neutral  & 55.7\% & 81.7\% & 66.2\% \\
\midrule
FinBERT & Positive & 75.9\% & 94.3\% & \textbf{84.1\%} \\
        & Negative & 82.5\% & 94.3\% & \textbf{88.0\%} \\
        & Neutral  & \textbf{97.0\%} & 53.3\% & 68.8\% \\
\bottomrule
\end{tabular}
\end{table}

\begin{table}[ht]
\centering
\caption{Confusion matrix for FinBERT (rows: true label; columns: predicted).}
\label{tab:cm}
\begin{tabular}{lccc}
\toprule
 & \textbf{Pred. positive} & \textbf{Pred. negative} & \textbf{Pred. neutral} \\
\midrule
True positive & \textbf{66} & 3 & 1 \\
True negative & 4 & \textbf{66} & 0 \\
True neutral  & 17 & 11 & \textbf{32} \\
\bottomrule
\end{tabular}
\end{table}

\begin{figure}[ht]
\centering
\includegraphics[width=0.85\textwidth]{metrics_barchart.pdf}
\caption{Per-class precision, recall and F1 for FinBERT (coloured bars) versus the
keyword baseline (grey bars) on the 200-sample benchmark. FinBERT's large recall
advantage on the negative class and the neutral asymmetry (high precision, low
recall) are immediately visible.}
\label{fig:bars}
\end{figure}

\begin{figure}[ht]
\centering
\includegraphics[width=0.6\textwidth]{confusion_matrix.pdf}
\caption{Row-normalised confusion matrix for FinBERT on the 200-sample benchmark.
The strong diagonal on the positive and negative classes contrasts with the
dispersed neutral row, visualising the neutral recall deficit.}
\label{fig:cm}
\end{figure}

\subsection{Explainability-guided diagnosis}
Applied to unambiguous cases, LIME assigns the highest weights to semantically
appropriate tokens (``surged'', ``raised'', ``guidance'' for a positive prediction
at 95\% confidence, Figure~\ref{fig:lime_pos}; ``decline'', ``sharp'', ``missing''
for a negative prediction at 97.6\%, Figure~\ref{fig:lime_neg}), confirming the
explanations are meaningful rather than artefactual. The diagnostically decisive
case is a misclassified neutral sentence: \emph{``The Federal Reserve kept interest
rates unchanged, maintaining its cautious outlook on inflation.''} FinBERT predicts
\emph{negative} (55.9\% confidence), and LIME identifies ``cautious'' and
``inflation'' as the driving negative weights, both legitimate risk signals in
isolation but contextually neutral here. This is precisely the error pattern
visible in Table~\ref{tab:cm}: the model conflates neutral-framed risk vocabulary
with genuine negative sentiment. LIME thus converts an aggregate symptom (53.3\%
neutral recall) into a specific, lexically-localised mechanism, exactly the kind of
actionable diagnosis that motivates a targeted intervention.

\begin{figure}[ht]
\centering
\includegraphics[width=0.8\textwidth]{lime_positive.pdf}
\caption{LIME token weights for a positive financial sentence. Green bars indicate
tokens supporting the \emph{positive} prediction; red bars indicate tokens pushing
against it. FinBERT predicted \emph{positive} at 95\% confidence.}
\label{fig:lime_pos}
\end{figure}

\begin{figure}[ht]
\centering
\includegraphics[width=0.8\textwidth]{lime_negative.pdf}
\caption{LIME token weights for a negative financial sentence. FinBERT predicted
\emph{negative} at 97.6\% confidence; \emph{decline}, \emph{sharp} and
\emph{missing} are the principal drivers.}
\label{fig:lime_neg}
\end{figure}

\subsection{Targeted mitigation}
Table~\ref{tab:finetune} reports the effect of fine-tuning on the disjoint
300-sample neutral-loaded corpus, re-evaluated on the unchanged benchmark. Neutral
recall improves by 18.4 percentage points (53.3\% $\rightarrow$ 71.7\%), the precise
effect the diagnosis predicted; neutral precision falls modestly (97.0\%
$\rightarrow$ 91.8\%), as expected when the model is encouraged to predict neutral
more often, but neutral F1 nonetheless rises by 11.7 points. Overall accuracy
increases to 86.5\% and macro F1 to 85.2\%, and polarised-class F1 is preserved (both
within 0.5 points of the stock model), confirming the intervention does not simply
trade one class error for another. The diagnosed failure mode is therefore
demonstrably addressable with a small, targeted dataset.

\begin{table}[ht]
\centering
\caption{Stock FinBERT versus neutral-focused fine-tuned FinBERT on the 200-sample
benchmark.}
\label{tab:finetune}
\begin{tabular}{lccccc}
\toprule
\textbf{Model} & \textbf{Accuracy} & \textbf{Macro F1} & \textbf{Neu.\ prec.} & \textbf{Neu.\ recall} & \textbf{Neu.\ F1} \\
\midrule
FinBERT (stock)      & 82.0\% & 80.3\% & \textbf{97.0\%} & 53.3\% & 68.8\% \\
FinBERT (fine-tuned) & \textbf{86.5\%} & \textbf{85.2\%} & 91.8\% & \textbf{71.7\%} & \textbf{80.5\%} \\
\midrule
\textit{Delta} & \textit{+4.5 pp} & \textit{+4.9 pp} & \textit{$-$5.2 pp} & \textit{+18.4 pp} & \textit{+11.7 pp} \\
\bottomrule
\end{tabular}
\end{table}

We are explicit about the strength of this evidence: the $+4.5$-point accuracy gain
is from a single 200-sample benchmark and a single fine-tuning seed, without
cross-validation, and the fine-tuned model has not been stress-tested on the full
live distribution. The result is reported as an indicative demonstration of the
diagnosis-to-mitigation loop, not as a definitive performance claim
(Section~\ref{sec:discussion}); a multi-seed cross-validated protocol is provided as
reproducible code (Section~\ref{sec:reprod}).

\subsection{Supporting analyses}
\paragraph{Sentiment-price correlation.}
A demonstrative run of the correlation engine across a nine-ticker panel (90-day
window, same-day alignment) produced modest, heterogeneous, and mostly
non-significant effects. Same-day Pearson coefficients ranged from $r = -0.10$ (SPY)
to $+0.24$ (NVDA); NVDA, AAPL and AMZN were individually significant at
$\alpha = 0.05$ (uncorrected), but under a Bonferroni correction for nine tickers
($\alpha_B = 0.0056$) none reached significance, and no Granger test attained
$p < 0.05$ at one- or two-day lags. SPY's near-zero coefficient is the expected null
for a diversified index tracker. These magnitudes are consistent with the literature
that sentiment-return effects are statistically detectable but economically small
\citep{tetlock2007,kraaijeveld2020}, and the rolling-correlation standard deviations
(0.14 to 0.34) show the relationship is unstable across sub-windows. We report this
as an honest null or modest result rather than selecting favourable windows.

\paragraph{Finance-emotion enrichment.}
Because no public benchmark annotates financial text at this emotion granularity, an
accuracy evaluation is not possible; we instead ran a component-contribution
ablation measuring how often removing a component flips the dominant emotion label on
the 200-sample benchmark. The domain lexicon is load-bearing (16.0\% of labels change
when it is removed; all samples collapse to ``mixed''), whereas zeroing the entropy
term flips no labels, an honest negative finding indicating the entropy component
requires re-tuning. The enrichment layer is therefore reported as deployed but only
partially validated.

\paragraph{System performance.}
FinBERT batch inference required 2.15~s for 200 samples ($\approx$11~ms per text);
the keyword baseline was effectively instantaneous; LIME explanations took 15 to
40~s, consistent with an on-demand rather than bulk feature.

% =====================================================================
\section{Discussion}\label{sec:discussion}
% =====================================================================
\subsection{Interpretation and significance}
The principal finding is that FinBERT's neutral-class deficit, usually reported as
an inert number, is a \emph{specific and repairable} failure mode. The conservatism
that produces 97.0\% neutral precision at 53.3\% recall is, in one sense, defensible
for a financial dashboard: mislabelling a strongly positive announcement as neutral
is arguably worse for a user than the reverse. But the 28/60 neutral error rate is a
non-trivial source of label noise, particularly for factual news. The contribution
of this paper is to show \emph{why} the errors occur and \emph{that} they can be
reduced: LIME localises the cause to neutral-framed risk vocabulary, and fine-tuning
on a small corpus of exactly that kind of text recovers 18.4 points of recall without
collateral damage to the polarised classes. This is a concrete instance of XAI used
as a diagnostic feedback mechanism rather than a presentation layer, the use case
that surveys identify as under-explored \citep{yeo2025}.

A second, methodological contribution is the demonstration that applied
financial-sentiment comparisons can and should be made with paired significance
testing. A 13.5-point accuracy gap is only a defensible claim once a McNemar test
($p = 0.002$) and non-overlapping bootstrap intervals rule out sampling variation;
many comparable studies report point estimates alone.

\subsection{Threats to validity and limitations}
We state limitations plainly. (i)~\emph{Benchmark scale and annotation.} The
benchmark is 200 samples, single-author-labelled, validated by an LLM cross-pass and
a single human annotator on a 100-sample subset; this is stronger than
single-annotator labelling but not equivalent to a multi-annotator panel of the kind
used to build Financial PhraseBank. Scaling the benchmark to several thousand samples
and commissioning a panel are priority future work. (ii)~\emph{Strength of the
mitigation evidence.} The fine-tuning result rests on a single seed and a single
benchmark without cross-validation, and the fine-tuned model has not been deployed or
evaluated on the full multi-source distribution; it evidences the
diagnosis-to-mitigation loop but is not a definitive performance claim.
(iii)~\emph{LIME stability.} LIME explanations are local linear approximations and
can vary across runs; they are corroborating evidence, not guarantees, and can in
principle be adversarially fooled \citep{slack2020}. (iv)~\emph{Data imbalance.} The
Twitter subcorpus is smaller and shallower than the Reddit and news subcorpora
because of X API v2 quota constraints, so Twitter is a secondary signal and
Twitter-only views carry lower statistical power. (v)~\emph{Correlation is not
causation.} The correlation engine is an exploratory tool; short-window sentiment
series are noisy and non-stationary, macro events confound sentiment and price
simultaneously, and the post-2008 liquidity regime weakens fundamental-to-price
linkages. The reported null Granger results should be read accordingly.
(vi)~\emph{Enrichment validation.} The finance-emotion layer is deployed but lacks an
accuracy evaluation against an annotated emotion corpus, which does not yet exist for
these classes.

\subsection{Ethical considerations}
All data are publicly available and anonymised at ingestion, in line with platform
terms of service and the institutional ethical framework. Two broader issues warrant
note. First, a sentiment-analytics tool confers an informational advantage on its
users, connecting to long-standing debates on market efficiency \citep{fama1970}.
Second, where natural-language summaries of sentiment are generated, model outputs
can present statistical artefacts as confident narratives; outputs should therefore
be accompanied by explicit documentation of intended use and limitations, an approach
consistent with the model-cards framework \citep{mitchell2019}. This study is for
academic purposes; commercial deployment would require additional governance.

% =====================================================================
\section{Conclusions}\label{sec:conclusion}
% =====================================================================
This paper reframed a familiar weakness of FinBERT, namely poor neutral-class recall,
as a diagnosable and repairable failure mode, and demonstrated an
explainability-guided loop that closes the gap between identifying a model error and
fixing it. On a purpose-built multi-source benchmark, FinBERT significantly
outperformed a finance-lexicon baseline (82.0\% vs 68.5\% accuracy; McNemar
$p = 0.002$), but its 53.3\% neutral recall was traced by LIME to a single mechanism:
contextually neutral risk vocabulary attributed to the negative class. A targeted
fine-tuning intervention on a small, disjoint corpus of neutral-but-lexically-loaded
text recovered neutral recall by 18.4 points and lifted accuracy to 86.5\% while
preserving polarised-class performance. The methodological message generalises beyond
FinBERT: local interpretability can be used as an instrument to localise systematic
errors precisely enough to motivate proportionate, measurable mitigations, and
applied sentiment comparisons should be reported with paired significance testing.

Future work follows directly from the limitations. The most valuable next steps are
(i)~scaling the benchmark and commissioning a multi-annotator panel; (ii)~validating
the fine-tuning result under cross-validation and on the full live distribution;
(iii)~constructing an annotated financial-emotion corpus to evaluate the enrichment
layer formally; and (iv)~extending the baseline comparison to current financial LLMs
\citep{wu2023,fatouros2023}. Each builds on the infrastructure and benchmark
established here.

\paragraph{Reproducibility.}\label{sec:reprod}
The evaluation harness, the paired-significance and bootstrap routines, a multi-seed
cross-validated fine-tuning protocol, and an external Financial PhraseBank
comparison are released as runnable code accompanying this submission (see the Data
availability statement).

% =====================================================================
%  DECLARATIONS
% =====================================================================
\section*{Use of Generative-AI tools declaration}
The author used a generative AI assistant (Anthropic Claude) in two ways. First,
as part of the research method, the model was used as an independent second
annotator to cross-validate the labels of the 200-sample evaluation benchmark
(Section~\ref{sec:benchmark}); it received only the raw texts and the labelling
rubric, with no access to the author's labels. Second, in preparing this
manuscript, the author used the same assistant to help edit and tighten the
author's own drafts. The author designed the study, produced and verified all data,
results,
analyses and conclusions, reviewed all generated text, and takes full
responsibility for the content. No AI tool is listed as an author. The Assistive-AI
language tools noted in the journal's instructions (grammar and spelling checkers)
are not generative tools and are not within the scope of this declaration.

\section*{Acknowledgments}
The author thanks Professor S.\ Lynch (Department of Computer Science, Loughborough
University) for supervision throughout the project, the finance and technology
practitioners who provided requirements input, and the volunteer annotator who
contributed to benchmark validation. This research did not receive any specific grant
from funding agencies in the public, commercial, or not-for-profit sectors. The
Article Processing Charge is covered by Guangzhou University.

\section*{Ethics statement}
The text analysed in this study was collected from publicly available sources
(Reddit, X/Twitter and online financial news) and was anonymised at the point of
ingestion, with author handles, display names and geolocation removed; no
personally identifying information is reported. Under the Loughborough University
ethical framework, this use of secondary, publicly available data did not require
formal ethical review. The wider project, including its human-participant component,
received a favourable ethical opinion from the Loughborough University Ethics Review
Sub-Committee through the University's ethical review system, with Professor S.\ Lynch
as responsible investigator. The volunteer who provided independent annotations for
benchmark validation took part with informed consent and contributed only sentiment
labels; no personal data were collected from that individual.

\section*{Conflict of interest}
The author declares no conflicts of interest in this paper.

\section*{Data availability}
The evaluation benchmark, derived results, and analysis code that support the
findings of this study are available from the author on reasonable request. The
underlying social-media and news text is subject to the respective platforms' terms
of service and cannot be redistributed in raw form.

% =====================================================================
%  REFERENCES  (inline; author-year via natbib bibitem labels)
%  Journal names abbreviated (ISO-4 style, no full stops, per the AIMS
%  example "J Amer Math Soc Volume: pages").
% =====================================================================
\begin{thebibliography}{99}
\setlength{\itemsep}{2pt}

\bibitem[Antweiler and Frank(2004)]{antweiler2004}
Antweiler W, Frank MZ (2004) Is all that talk just noise? The information content of
internet stock message boards. \emph{J Finance} 59: 1259--1294. \url{https://doi.org/10.1111/j.1540-6261.2004.00662.x}

\bibitem[Araci(2019)]{araci2019}
Araci D (2019) FinBERT: Financial sentiment analysis with pre-trained language
models. \emph{arXiv:1908.10063}. \url{https://doi.org/10.48550/arXiv.1908.10063}

\bibitem[Baker and Wurgler(2006)]{baker2006}
Baker M, Wurgler J (2006) Investor sentiment and the cross-section of stock returns.
\emph{J Finance} 61: 1645--1680. \url{https://doi.org/10.1111/j.1540-6261.2006.00885.x}

\bibitem[Bollen et~al.(2011)]{bollen2011}
Bollen J, Mao H, Zeng X (2011) Twitter mood predicts the stock market. \emph{J Comput
Sci} 2: 1--8. \url{https://doi.org/10.1016/j.jocs.2010.12.007}

\bibitem[Cohen(1960)]{cohen1960}
Cohen J (1960) A coefficient of agreement for nominal scales. \emph{Educ Psychol
Meas} 20: 37--46. \url{https://doi.org/10.1177/001316446002000104}

\bibitem[Da et~al.(2015)]{da2015}
Da Z, Engelberg J, Gao P (2015) The sum of all FEARS investor sentiment and asset
prices. \emph{Rev Financ Stud} 28: 1--32. \url{https://doi.org/10.1093/rfs/hhu072}

\bibitem[Deng et~al.(2023)]{deng2023}
Deng X, Bashlovkina V, Han F, et~al. (2023) LLMs to the moon? Reddit market sentiment
analysis with large language models. \emph{arXiv:2301.10716}. \url{https://doi.org/10.1145/3543873.3587605}

\bibitem[Devlin et~al.(2019)]{devlin2019}
Devlin J, Chang MW, Lee K, et~al. (2019) BERT: Pre-training of deep bidirectional
transformers for language understanding. \emph{Proc NAACL-HLT 2019}, 4171--4186. \url{https://doi.org/10.18653/v1/N19-1423}

\bibitem[Dickey and Fuller(1979)]{dickey1979}
Dickey DA, Fuller WA (1979) Distribution of the estimators for autoregressive time
series with a unit root. \emph{J Am Stat Assoc} 74: 427--431. \url{https://doi.org/10.2307/2286348}

\bibitem[Dietterich(1998)]{dietterich1998}
Dietterich TG (1998) Approximate statistical tests for comparing supervised
classification learning algorithms. \emph{Neural Comput} 10: 1895--1923. \url{https://doi.org/10.1162/089976698300017197}

\bibitem[Doshi-Velez and Kim(2017)]{doshivelez2017}
Doshi-Velez F, Kim B (2017) Towards a rigorous science of interpretable machine
learning. \emph{arXiv:1702.08608}. \url{https://doi.org/10.48550/arXiv.1702.08608}

\bibitem[Efron and Tibshirani(1993)]{efron1993}
Efron B, Tibshirani RJ (1993) \emph{An Introduction to the Bootstrap}, New York:
Chapman and Hall/CRC. \url{https://doi.org/10.1201/9780429246593}

\bibitem[Fama(1970)]{fama1970}
Fama EF (1970) Efficient capital markets: A review of theory and empirical work.
\emph{J Finance} 25: 383--417. \url{https://doi.org/10.2307/2325486}

\bibitem[Fatouros et~al.(2023)]{fatouros2023}
Fatouros G, Soldatos J, Kouroumali K, et~al. (2023) Transforming sentiment analysis
in the financial domain with ChatGPT. \emph{Mach Learn Appl} 14: 100508. \url{https://doi.org/10.1016/j.mlwa.2023.100508}

\bibitem[Fleiss(1971)]{fleiss1971}
Fleiss JL (1971) Measuring nominal scale agreement among many raters. \emph{Psychol
Bull} 76: 378--382. \url{https://doi.org/10.1037/h0031619}

\bibitem[Garcia(2013)]{garcia2013}
Garcia D (2013) Sentiment during recessions. \emph{J Finance} 68: 1267--1300. \url{https://doi.org/10.1111/jofi.12027}

\bibitem[Granger(1969)]{granger1969}
Granger CWJ (1969) Investigating causal relations by econometric models and
cross-spectral methods. \emph{Econometrica} 37: 424--438. \url{https://doi.org/10.2307/1912791}

\bibitem[Huang et~al.(2023)]{huang2023}
Huang AH, Wang H, Yang Y (2023) FinBERT: A large language model for extracting
information from financial text. \emph{Contemp Account Res} 40: 806--841. \url{https://doi.org/10.1111/1911-3846.12832}

\bibitem[Kraaijeveld and De~Smedt(2020)]{kraaijeveld2020}
Kraaijeveld O, De~Smedt J (2020) The predictive power of public Twitter sentiment for
forecasting cryptocurrency prices. \emph{J Int Financ Mark Inst Money} 65: 101188. \url{https://doi.org/10.1016/j.intfin.2020.101188}

\bibitem[Landis and Koch(1977)]{landis1977}
Landis JR, Koch GG (1977) The measurement of observer agreement for categorical data.
\emph{Biometrics} 33: 159--174. \url{https://doi.org/10.2307/2529310}

\bibitem[Leippold(2023)]{leippold2023}
Leippold M (2023) Sentiment spin: Attacking financial sentiment with GPT-3.
\emph{Finance Res Lett} 55: 103957. \url{https://doi.org/10.1016/j.frl.2023.103957}

\bibitem[Long et~al.(2023)]{long2023}
Long S, Lucey B, Xie Y, et~al. (2023) ``I just like the stock'': The role of Reddit
sentiment in the GameStop share rally. \emph{Financ Rev} 58: 19--37. \url{https://doi.org/10.1111/fire.12328}

\bibitem[Lopez-Lira and Tang(2023)]{lopezlira2023}
Lopez-Lira A, Tang Y (2023) Can ChatGPT forecast stock price movements? Return
predictability and large language models. \emph{SSRN Working Paper 4412788}. \url{https://doi.org/10.2139/ssrn.4412788}

\bibitem[Loughran and McDonald(2011)]{loughran2011}
Loughran T, McDonald B (2011) When is a liability not a liability? Textual analysis,
dictionaries, and 10-Ks. \emph{J Finance} 66: 35--65. \url{https://doi.org/10.1111/j.1540-6261.2010.01625.x}

\bibitem[Lundberg and Lee(2017)]{lundberg2017}
Lundberg SM, Lee SI (2017) A unified approach to interpreting model predictions.
\emph{Adv Neural Inf Process Syst} 30: 4765--4774. \url{https://doi.org/10.48550/arXiv.1705.07874}

\bibitem[Malo et~al.(2014)]{malo2014}
Malo P, Sinha A, Takala P, et~al. (2014) Good debt or bad debt: Detecting semantic
orientations in economic texts. \emph{J Assoc Inf Sci Technol} 65: 782--796. \url{https://doi.org/10.1002/asi.23062}

\bibitem[Mitchell et~al.(2019)]{mitchell2019}
Mitchell M, Wu S, Zaldivar A, et~al. (2019) Model cards for model reporting.
\emph{Proc Conf Fairness Account Transpar (FAT* 2019)}, 220--229. \url{https://doi.org/10.1145/3287560.3287596}

\bibitem[Pontiki et~al.(2014)]{pontiki2014}
Pontiki M, Galanis D, Pavlopoulos J, et~al. (2014) SemEval-2014 Task 4: Aspect based
sentiment analysis. \emph{Proc 8th Int Workshop Semant Eval (SemEval 2014)}, 27--35. \url{https://doi.org/10.3115/v1/S14-2004}

\bibitem[Ribeiro et~al.(2016)]{ribeiro2016}
Ribeiro MT, Singh S, Guestrin C (2016) ``Why should I trust you?'': Explaining the
predictions of any classifier. \emph{Proc 22nd ACM SIGKDD Int Conf Knowl Discov Data
Min}, 1135--1144. \url{https://doi.org/10.1145/2939672.2939778}

\bibitem[Rudin(2019)]{rudin2019}
Rudin C (2019) Stop explaining black box machine learning models for high stakes
decisions and use interpretable models instead. \emph{Nat Mach Intell} 1: 206--215. \url{https://doi.org/10.1038/s42256-019-0048-x}

\bibitem[Schumaker and Chen(2009)]{schumaker2009}
Schumaker RP, Chen H (2009) Textual analysis of stock market prediction using
breaking financial news: The AZFinText system. \emph{ACM Trans Inf Syst} 27:
12:1--12:19. \url{https://doi.org/10.1145/1462198.1462204}

\bibitem[Slack et~al.(2020)]{slack2020}
Slack D, Hilgard S, Jia E, et~al. (2020) Fooling LIME and SHAP: Adversarial attacks
on post hoc explanation methods. \emph{Proc AAAI/ACM Conf AI Ethics Soc (AIES 2020)},
180--186. \url{https://doi.org/10.1145/3375627.3375830}

\bibitem[Sravankumar et~al.(2025)]{sravankumar2025}
Sravankumar B, Krishna G, Rao M (2025) Integration of explainable AI with deep
learning for real-time sentiment-driven stock price prediction. \emph{Comput Econ}. \url{https://doi.org/10.1007/s10614-025-11159-w}

\bibitem[Tetlock(2007)]{tetlock2007}
Tetlock PC (2007) Giving content to investor sentiment: The role of media in the
stock market. \emph{J Finance} 62: 1139--1168. \url{https://doi.org/10.1111/j.1540-6261.2007.01232.x}

\bibitem[Vaswani et~al.(2017)]{vaswani2017}
Vaswani A, Shazeer N, Parmar N, et~al. (2017) Attention is all you need. \emph{Adv
Neural Inf Process Syst} 30: 6000--6010. \url{https://doi.org/10.48550/arXiv.1706.03762}

\bibitem[Wu et~al.(2023)]{wu2023}
Wu S, Irsoy O, Lu S, et~al. (2023) BloombergGPT: A large language model for finance.
\emph{arXiv:2303.17564}. \url{https://doi.org/10.48550/arXiv.2303.17564}

\bibitem[Yang et~al.(2020)]{yang2020}
Yang Y, Uy MCS, Huang A (2020) FinBERT: A pretrained language model for financial
communications. \emph{Proc IJCAI 2020 FinNLP Workshop}. \emph{arXiv:2006.08097}. \url{https://doi.org/10.48550/arXiv.2006.08097}

\bibitem[Yeo et~al.(2025)]{yeo2025}
Yeo WJ, van~der~Heever W, Mao R, et~al. (2025) A comprehensive review on financial
explainable AI. \emph{Artif Intell Rev} 58. \url{https://doi.org/10.1007/s10462-024-11077-7}

\end{thebibliography}

\end{document}
